% Preparatory outline only. No production manuscript or theorem numbering changes.
% Author of this outline: /root/astra_quartic_opt05_anchor, 2026-11-07.
% Exact source and review bindings: APPLICATION-DELTA.json and INTEGRATION-MAP.json.
% The proposed labels below are local to this draft and are not accepted exposition.
\documentclass{article}
\usepackage{amsmath,amssymb}
\begin{document}
\paragraph{Revision 4 scope.} The inherited sections below retain the revision 3 cutoff. The appended revision 4 section supplies the selected 313-variable route; old 315-specific scope statements are historical, not current approval restrictions. This remains a preparatory outline.

\title{Proof integration outline for the accepted 315-variable construction}
\author{Preparatory notes}
\date{November 7, 2026}
\maketitle

\section{Scope and fixed conventions}\label{outline:scope}
The selected object is \texttt{combined315}, identified by its full IR hash
in \texttt{CANDIDATES.json}. A variable count is not a candidate identifier.
This document plans the proof exposition; it is neither a new proof audit
nor a release manuscript. The strict publication threshold of fewer than
300 total variables is unmet.

Keep the pinned Version 13 source editions, errata qualifications, restricted
converses, field basis and valuation normalizations. Write
\[
 K=\mathbb Q(\rho),\quad \rho^2=2,\quad D=(3-\rho)^2,
 \quad C=3^{2^{25}}.
\]
The scalar coefficient is $F_s=2g_0g_1$. The root abscissa coefficients
are $X_A=2v_{83}$, $X_T=2v_{87}$ and $X_Z=2v_{89}$.
These are different functions. Do not replace $F_s$ by $g_1^2$.
All notation for point indices, contact primes and scales must first be
identified with the retained Version 13 definitions.

\section{Exact chart changes and their exceptional branches}
\label{outline:charts}
\paragraph{Entries N01 and N02.}
State the exact Kummer affine substitution and its reverse with all named
divisors. Then state squared-H reconstruction only on $UQ\ne0$.
For the latter, the recovered quantities are
\[
 f=(U+W)/Q,\qquad F=(V+W)/Q.
\]
Give the literal identities proving the former curve and subtraction
equations, including each physical affine normalization. The source
reports bound to N01 and N02 supply these identities. The construction
passage must name every altered scalar, primary and logarithm consumer.

\paragraph{Entry N03.}
Make a separate lemma for the newly admitted $U=0$ and $Q=0$ cases.
It must classify their fixed values, prove they are units at every
consumed prime, and obtain contact radical one at all three scales.
Use the restored scalar anchor and root-height inequality to exclude
them before giving H a group-index interpretation. No logarithm ratio
or H-dependent denominator may be used to prove this exclusion.
The precise branch constants and their source spans are bound in N03;
they must be transcribed with the source's meanings of $U,V,v,c$.

\section{Contact parameter and conditional guards}
\label{outline:contact}
\paragraph{Entries N04, N05 and N06.}
Replace the former weighted contact parameter by
\[
 R=F_5(A\lambda),\qquad E=R^{-1},\qquad
 B=\frac{R}{T^2(1+R)^2}.
\]
Here $A$ is the detector scale, not the unrelated Lucas field input
used later. State soundness at a required positive odd valuation
$\alpha=v_p(R)$, where $1+R$ is a unit. Prove every retained pole and
norm valuation using the actual square denominator.

The guard-removal statement is conditional at required contact primes.
Do not promote it to a global trace inverse when a contact radical is one.
The completeness lemma must separately order finite support, small-prime
depths, initial real choices, root choices, compact neighborhoods and final
CRT progression. Quote the original open-set and density hypotheses,
not independent local and real choices made after the same point is fixed.

\section{Root parity and the shifted quaternion application}
\label{outline:anchor}
\paragraph{Entries N15, N16 and N17.}
First obtain integral root indices from the retained square-abscissa charts.
Do not assume oddness. State and prove the coefficient formulas used by
the 359 component at both 3 and 2. In particular, the 3-adic valuation
of $X_n$ is zero for odd $n$, while for nonzero even $n$ it is
$-2-2v_3(n)$; at 2 an odd root has valuation $-1$.
State separately every zero-index exclusion supplied by the chart.

Define
\[
 q=12(X_A^2+X_T^2)+108X_Z^2
   =48v_{83}^2+48v_{87}^2+432v_{89}^2.
\]
The changed application of the inherited quaternion criterion has
parameter $F_s/C+q$. Its source must cover zero as well as nonzero
parameters. The three competing negative valuation types lie in
$2\mathbb Z$, $1+4\mathbb Z$ and $3+4\mathbb Z$.
Use the unique-minimum argument to recover root oddness and the old
scalar anchor. At 2 use the integral formal subgroup, without claiming
good reduction there.

Then give the three full A,T,Z sign actions. Check them on the final
detector equations, not only on the original point blocks. They restore
the old half-charts while fixing actual root points, root abscissae,
detector outputs and the distinguished input.

\paragraph{Entries N07 and N08.}
Only after restoring the 326 source may its inherited 366 argument prove
S oddness and apply the S sign normalization and half-point restoration.
The S action is not asserted directly on the shifted private B equation.
Retain the full coordinate lists and the existential half-point argument
from their pinned source and review passages.

\section{Valuations before H and the three common-place groups}
\label{outline:detectors}
\paragraph{Entries N09, N10, N13 and N14.}
State the unconditional Kummer-to-curve and primary fourth-power
identities before interpreting H. Prove the generic valuation classes
from those identities and handle the fixed-unit exceptional leaves
separately. Record the earlier 354 reviewer's valuation-only repair
as an input to this later application.

For the selected three-group construction, use exactly
\begin{align*}
 G_0&=1+E^{158}H_1^4+E^{157}H_4^4+E^{39}M_0^4+E^8T_0+E^4T_4,\\
 G_1&=1+E^{158}H_2^4+E^{79}H_0^2+E^8T_2+E^4Z_1+EP_1,\\
 G_2&=1+E^{158}H_3^4+E^{39}M_1^4+E^8T_3+E^{157}P_2^2+E^4P_3^2+EP_4^2.
\end{align*}
Bind every letter to its actual cached field or rational expression.
For each group, tabulate all possible negative valuations modulo 16,
including the actual powers and weights. Explain why the minimum cannot
cancel. A zero leaf has infinite valuation and must be handled explicitly.
Then derive all leaf conditions at the same place selected by the
combined norm argument; separately successful places are insufficient.

For completeness retain the original $\Pi^{19}$ progression, primary
formal depth 20 and reciprocal depth 40. The shallow G2 ratio has actual
precision 12, not merely the required order 10. The other high ratios
have precision 40 and primary comparisons have precision 80.
The actual positive-margin lists, in the source term order, are
\[
 (2,3,9,8,12),\qquad (2,1,8,4,79),\qquad
 (2,281,8,3,156,159).
\]
In particular, the $E^{157}P_2^2$ term consumes the stronger U2 depth.
Do not import the separate strong-H330 branch's $\Pi^{39}$ choice into
this proof. N11 documents that unselected branch, not an extra premise
of the selected construction.

\section{Exact private anchor unscaling}
\label{outline:unscale}
\paragraph{Entry N18.}
State a full solution-set equivalence under the retained guard
$g_0\tau=1$. Starting from
\[
 Q=q,\quad B=2g_1+CQ\tau,\quad
 U^2+6V^2-2W^2-4C\tau W-3B^2=0,
\]
put $u=g_0U$, $v=g_0V$, $w=g_0W$, $T=g_0B$. The new equations are
\[
 T=2g_0g_1+Cq,\qquad u^2+6v^2-2w^2-4Cw-3T^2=0.
\]
The inverse sets $Q=q$, $B=\tau T$, $U=\tau u$, $V=\tau v$,
$W=\tau w$. Both maps fix every nonprivate coordinate, and the guard
makes their compositions the identity. The old anchor is $\tau^2$
times the new anchor. Include all five unconditional residual identities
from the accepted exact319 map to certify the literal row normalizations.
No private witness is assumed nonzero. This changes no quaternion parameter.

In the three-group composition the anchor row labels 903 and 904 become
1903 and 1904. Only row labels change. The detector rows 903 and 904
remain present. Unscaling deletes anchor row 1904 and coordinate 904,
retaining coordinate 903 as T. Keep these namespaces explicit.

\section{The shorter exact Laurent chain}
\label{outline:chain}
\paragraph{Entries N12 and N20.}
For each detector retain $ER=1$ and write the new equations as
\begin{align*}
 EJ&=R^2,&E_4J&=E,&E_8&=E_4^2,&E_{19}J&=E_8^2,\\
 RE_{39}&=E_{19}^2,&RE_{79}&=E_{39}^2,&
 E_{158}&=E_{79}^2,&E_{157}&=RE_{158}.
\end{align*}
Since $ER=1$, both factors are nonzero before any arithmetic
interpretation. The first equation gives $J=R^3=E^{-3}\ne0$.
Successive substitution therefore gives $E_i=E^i$ for every displayed
index. All squarings assign new values; none is inverted by choosing a root.

The forward map deletes the old $E_2$ and replaces the old $E_{11}$
helper by $R^3$. The inverse restores $E_2=E^2$ and $E_{11}=E^{11}$.
Both maps fix all consumed outputs and all other coordinates, including
the private anchor. Hence the common-place and completeness arguments
receive exactly their old inputs. The literal three copies delete
coordinates 361,409,457 and rows 313,363,413; they repurpose coordinates
355,403,451. Only rows 544,546,550,552,556,558,745,760,775 change.
The frozen 315 composition checks all these consumers and the three root
sign actions again. No new local condition is introduced.

\section{Order of the complete argument}
\label{outline:composition}
For soundness, invert the chain map to obtain 318, invert private
unscaling to obtain the three-anchor319 system, prove root parity and
restore the 326 source, then apply the inherited S normalization.
Next prove pre-H valuation claims and exclude the exceptional branches.
Only then use the retained logarithm, integrality-transfer, scalar24,
height and input arguments.

For completeness, construct the 326 witnesses with its ordered
$\Pi^{19}$ finite and real choices. Apply the 359 private quaternion
extension after roots and scalar data are fixed. Apply exact unscaling,
then assign the unique shortened-chain helpers. These last steps leave
all earlier choices unchanged.

\section{Separately accepted variants}
\label{outline:variants}
\paragraph{Entry N19, unselected.}
The squared-$g_1$ alternative has parameter $g_1^2/C+q$ and requires its
own orbit-specific valuation and 2-integrality lemmas. It is not used in
315 and must not be combined additively with N18.

\paragraph{Entry N21, only on exact319.}
Add $Q=q$ and an output-linear weighted constant chain
\[
 c_0=3^{a_0},\qquad c_i=3^{a_i}c_{i-1}^2.
\]
The recurrence $e_i=2e_{i-1}+a_i$ forces unique rational values
$c_i=3^{e_i}$ and final exponent $2^{25}$. Replace fixed C in the two
anchor consumers by $c_{\rm last}$ and write $T=2g_0g_1+c_{\rm last}Q$.
The separate Q gate is essential for quadratic degree in this lift.
Forgetting the unique new values is the inverse. Record three different
IRs, with 334,335 and336 variables, rather than calling them one ordinary319.
Their exact maximum coefficients include the inherited constant K:
$3^{4098}+K$, $3^{2050}+K$ and $3^{1026}+K$ respectively.
These accepted artifacts have not been transported to 315 by this outline.

\paragraph{Entry N22, only on combined320.}
For the Lucas field inputs, rename the old private product $\mathcal C=AZ$
as $S=\mathcal C-2Z$. With $V=AT$ the identity is
\[
 F+1=T^2+4Z^2-4Z\mathcal C+2ZV+\mathcal C^2-\mathcal CV
     =T^2+S^2-SV.
\]
The inverse adds $2Z$; no divisor or exceptional case occurs. Retain
physical output components $v_{253}$ and $2v_{254}$. Only coordinates
249,250 and their four rows change. This saves 61 monomials and no
variables on the reviewed source. A transport to 315 needs a separate
full consumer and coefficient check.

\paragraph{Entry N23, optional appendix only.}
The bound $A(2^n-1)$ applies only to the stated isolated acyclic monomial
constant-chain model with counted seeds and the separate Q gate. Its
coefficient-cap argument must retain those restrictions. It supplies no
unrestricted optimality result and is not a premise of integer projection.

\section{Finite encoding and inventory closure}
\label{outline:finite}
For the frozen selected 315 IR, distinguish 5732 coefficient-power entries
from 5729 ordinary monomials after fixing C. The three collisions are
at the fourth powers of coordinates 83,87,89. The exact maximum is
$186624C^2+64$. Preserve the complete primitive-residual SOS certificate,
active-coordinate count, degree four and integer-content-one checks.
The rational SOS equivalence must be stated separately from the arithmetic
integer-defining argument.

Before any manuscript integration, write full application contracts for
each changed use, with the original source qualifications and all nine
required obligations. Bind construction and consumer spans to the final
text. The machine-readable integration map assigns each N-entry to the
specific Version 13 applications and source consumer spans. Acceptance
of the component packets does not approve this draft exposition.
\section{Revision 4: common scale and constrained anchor witnesses}
\subsection{N24-COMMON-SCALE-LOCAL}
\label{outline:r4-n24-common-scale-local}
At one successful place alpha=v(R)>0 is odd. Residue classes are distinct for every integer c=v(gamma), before H is interpreted. The guard forces c>15alpha/32. The same place yields high order40, middle order10, primary f(U2/U3/U4) depths40/20/22 and all low formality conditions. No other R is assumed integral. Preserve pre-H signed branches and empty-radical height exclusion.
\par Inherited pre-H Kummer periods and trace/pole theorem. Distinct residue classes and lower bounds are proved internally; no new external theorem.
\par Consumers: Shared314 rows900..917 rescale H/M/P outputs;2000..2009 define gamma and P1 square; group rows918..950 and unchanged W/norm/trace/pole rows. Pre-H branch exclusions precede all logarithm use.
\subsection{N25-H-DEPTH-BOOTSTRAP}
\label{outline:r4-n25-h-depth-bootstrap}
At depth one use log(Q)=z(Q) times a unit, extended by1 at Q=O. One of half-indices h,h-1 is a unit. Positive cross-tests force both ratios integral. Since F(0,0,0)=-1, at least one Li is a unit. Exact z(Ui)=Li*z(H) and primary depth at least10 force H depth10. Only then use relative primary precision40 and alternate precision12/40 to derive integer-label congruences. F on integer labels is not assumed during bootstrap.
\par Retained formal logarithm unit identity, identity extension and exact primary charts. The bootstrap is a new internal lemma, not a new external arithmetic assumption.
\par Consumers: Exact primary I*L=V*S\_H equations; physical F+1=(v253,2v254), F\textasciicircum{}2=(v319-2v253+1,4v320-4v254); rows906/907 send F\textasciicircum{}4/gamma160 to H4. Source v13 lines20885..20900 gives the depth-one unit identity.
\subsection{N26-COMMON-SCALE-COMPLETE}
\label{outline:r4-n26-common-scale-complete}
Fix finite neighborhoods and real bounds first. All eight new common coefficients have positive small-prime slopes155,150,145,35,30,65,100,24 in nu and negative real exponents. Select contacts, gamma and final H in the stated order. Repeated contacts impose one valuation. Completeness uses actual G2 precision12 and primary depth20, not a strengthened Pi progression. Additional gamma denominator primes do not alter pole parity away from odd R-support.
\par Retained simultaneous-prime theorem and real-progression lemma. CRT prescription, coefficient estimates and parity calculation are internal.
\par Consumers: All three G0/G1/G2 evaluations and unchanged scalar/pole witnesses; same final H simultaneously satisfies finite and real conditions.
\subsection{N27-SPARSE-SHARED314}
\label{outline:r4-n27-sparse-shared314}
Subtract/add2Z are inverse polynomial maps on every rational assignment. Only rows201,202,205,206 change. The entire F/F+1 interface, all gamma/group consumers, and every zero case are fixed.
\par Polynomial identity over Q(rho), rho\textasciicircum{}2=2. No external theorem or new completeness choice.
\par Consumers: Rows201/202 and205/206; unchanged203/204 define V; output253/254 consumed by271/272 and906/907. Exactly6074 monomials and314 variables.
\subsection{N28-ORDINARY-SHARED314}
\label{outline:r4-n28-ordinary-shared314}
Every gate is output-linear and forces a unique rational value. Exponent recurrence reaches2\textasciicircum{}25. Q=48v83\textasciicircum{}2+48v87\textasciicircum{}2+432v89\textasciicircum{}2. Only rows504/1903 change. Forgetting added coordinates is inverse to unique extension. No gamma restriction. Maxima3\textasciicircum{}1026+K,3\textasciicircum{}2050+K,3\textasciicircum{}4098+K have1627/3250/6496 bits; 330 exceeds3247-bit cap.
\par Exact rational algebra, integer exponent induction and full finite tables only.
\par Consumers: New rows/coordinates1000..1016 or shorter; anchor504/1903 only. Detector903/904 remain unchanged. No ordinary lift of anchorGamma313 is included.
\subsection{N29-ANCHOR-GAMMA-CHAIN}
\label{outline:r4-n29-anchor-gamma-chain}
u16*J=1 forces u nonzero. Every helper is uniquely determined. Reconstruct old gamma and powers to recover every source row before interpreting points or valuations. This proves soundness only; arbitrary old witness fibers are not claimed equivalent.
\par Laurent reconstruction under an explicit inverse. No new external theorem; new forward existence is N30..N33.
\par Consumers: Exactly rows919,930,941,2000,2005 change. Row504 anchor and1903 T remain fixed; inverse2004 retained. All278 other residuals unchanged.
\subsection{N30-ANCHOR-SMALL-LOCAL}
\label{outline:r4-n30-anchor-small-local}
Before scalar selection impose Fs in8Z2 using permitted further formal depths. Odd-root X\_i valuation-1 implies q=1mod8 and t=T/C odd. At2 use v=T/2,z=C/2 sqrt(4-3t\textasciicircum{}2+2(u/C)\textasciicircum{}2). At3 use v=0,z=C sqrt(1-3t\textasciicircum{}2/2+(u/C)\textasciicircum{}2/2), nu3>2\textasciicircum{}25. At17 use v=T/sqrt2,z=C sqrt(1+u\textasciicircum{}2/(2C\textasciicircum{}2)). No integrality of T at17 is assumed.
\par Retained formal-subgroup finite-index/real-progression application and odd-root coefficient lemma. The local formulas are proved internally. The usual Q2 and odd-Qp unit-square criteria are standard external lemmas invoked without reproving them; no new global arithmetic theorem is introduced.
\par Consumers: Unchanged504/1903 define fixed T=Fs+Cq. Scalar513/514 contain no final H. v13 lines21430..21439 permit extra scalar depths;359 odd-abscissa lemma supplies exact X\_i convention.
\subsection{N31-ANCHOR-CONTACT-LOCAL}
\label{outline:r4-n31-anchor-contact-local}
Exclude denominators of fixed T before choosing primes. At p>76176 choose u=p\textasciicircum{}-1 and solve6V\textasciicircum{}2-2Z\textasciicircum{}2=p\textasciicircum{}2N-1. The nonzero residue is represented by a split product or finite-field norm, and a nonzero derivative gives a one-variable lift. This does not assume the other R values integral.
\par Same simultaneous-prime theorem with enlarged fixed exclusions. Norm surjectivity is deduced internally from the standard cyclicity theorem for finite-field multiplicative groups. Cyclicity and one-variable Hensel lifting are invoked standard external lemmas, not proved in the note and not stronger prime-pattern assumptions.
\par Consumers: Finite set of selected contacts for all three scales; fixes local open targets only, before one global anchor choice.
\subsection{N32-ANCHOR-RATIONAL-APPROXIMATION}
\label{outline:r4-n32-anchor-rational-approximation}
N is nonzero. Starting from a rational point P, Phi(d)=P-2B(P,d)d/Q(d) covers the complement of the tangent section. Explicit square-root tangent curves prove this complement meets every local target open set. Approximate three direction coordinates by CRT and a growing denominator1modM, then apply continuous Phi. Direction parameters are proof tools, not quartic variables. No general weak approximation theorem for quadrics is assumed.
\par Reuses the same quaternion representation theorem once for a basepoint. Reflection identity, local density reduction and coordinatewise approximation construction are proved internally; standard CRT and the local square-root-near1 lemma are invoked without reproof. No external weak-approximation theorem for quadrics is invoked.
\par Consumers: Creates existing private coordinates55,56,57; gamma derives from55. Rational basepoint comes from old quaternion criterion on t integral at2,3.
\subsection{N33-ANCHOR-COMPLETE-ORDER}
\label{outline:r4-n33-anchor-complete-order}
Fix scalar/root/T/basepoint; finite neighborhoods and bounds; denominator depths; contacts avoiding fixed T denominators; lambda; anchor approximation; final H in original Pi19 progression. T and test bounds are independent of private anchor/final H. Gamma has exactly the old uniform real and finite conditions. No new H-index congruence or prime-support restriction is introduced. Extra denominator primes preserve pole parity.
\par Same prime-selection, real-progression, preparation-zero and pole results with newly checked inputs. No self-approval of the contributed proof.
\par Consumers: All three detector/pole evaluations and final scalar preparation use one final H; unchanged root/scalar data determine1903 before approximation.
\end{document}
